\section{Marco teórico y trabajos relacionados}
Los sistemas de rastreo en tiempo real han mostrado impacto tanto en seguridad ciudadana como en logística: soluciones antirrobo basadas en GPS-GSM y plataformas de distribución apoyadas en la nube resaltan la importancia de la actualización continua de ubicaciones y la confiabilidad de los datos \cite{gorret2025vehicle, logistica2025seguimiento}.

La arquitectura es crítica cuando crecen los dispositivos o la frecuencia de envío. El caso Kbus documenta cómo la transición de un monolito a microservicios permitió sostener millones de eventos diarios sin degradar tiempos de respuesta, ofreciendo lecciones sobre desacoplamiento y escalabilidad \cite{torres2020kbus}.

En el cliente móvil, comparativas entre Flutter y React Native muestran que la elección depende de la experiencia del equipo y de la disponibilidad de librerías de mapas, siendo React Native atractivo por su ecosistema JavaScript \cite{macias2021flutterrn}. Experiencias previas de geolocalización en este framework sirven como guía para integrar rutas y posiciones \cite{dedios2021geolocalizacion}.

Para la mensajería, MQTT destaca por su patrón pub-sub ligero y la capacidad de ajustar la calidad de servicio, lo que lo convierte en la opción habitual en proyectos móviles con recursos limitados \cite{ortiz2022smolcomm, terrones2022presencia}.

La seguridad y la autenticación multiplataforma requieren gestionar credenciales de forma consistente y proteger los canales de comunicación; existen propuestas de autenticación unificada para React/React Native y buenas prácticas de programación segura que respaldan estas decisiones \cite{ye2022autenticacion, madrigal2020seguridad}. Finalmente, la privacidad sigue siendo un eje transversal en aplicaciones de rastreo y debe considerarse desde el diseño \cite{react2021}.
